\section{工程 Checklist：写分布式训练代码前先问什么}

\subsection{状态在哪里}

第一组问题是状态放置：

\begin{itemize}
    \item 参数、梯度、optimizer state、activation 是否复制？
    \item 哪些对象被 sharding？按 batch、width、depth、sequence 还是 expert？
    \item 哪个时刻需要完整副本？这个副本是长期保存还是临时 all-gather？
\end{itemize}

\subsection{通信何时发生}

第二组问题是通信时序：

\begin{itemize}
    \item collective 是 all-reduce、reduce-scatter、all-gather 还是 all-to-all？
    \item 通信发生在每层、每个 block、每个 microbatch，还是每个 step？
    \item 是否能与 backward compute overlap？
    \item 是否有 barrier 或隐式同步破坏 overlap？
\end{itemize}

\subsection{硬件是否匹配策略}

第三组问题是拓扑匹配：

\begin{itemize}
    \item 高频通信是否留在 NVLink/NVSwitch 域内？
    \item 跨节点链路是否只承载较粗粒度通信？
    \item NCCL benchmark 的实际带宽是否支持理论计划？
    \item 最慢 rank 或最慢 stage 在哪里？
\end{itemize}

\begin{importantbox}{最终判断标准}
一个并行策略的好坏不由名字决定，而由每个 step 的端到端吞吐、显存峰值、扩展效率和稳定性决定。可解释的系统账本比“用了某某 parallelism”更重要。
\end{importantbox}

\subsection{本章小结}

写分布式训练代码前，先画状态图和通信图。状态图告诉你什么被复制、什么被分片；通信图告诉你每个同步点移动什么张量、走什么链路。没有这两张图，调参只是猜。

